\documentclass[11pt]{article}
\usepackage[margin=0.65in]{geometry}
\usepackage{xcolor}
\usepackage{tcolorbox}
\usepackage{booktabs}
\usepackage{array}
\usepackage{amsmath}
\usepackage{fancyhdr}
\usepackage{listings}
\usepackage{enumitem}
\usepackage{hyperref}
\usepackage{multicol}
\usepackage{graphicx}
\usepackage{titlesec}

\definecolor{DeepNavy}{HTML}{0B1F3A}
\definecolor{LightBlue}{HTML}{EAF3FF}
\definecolor{LightGray}{HTML}{F4F6F8}
\definecolor{CodeBg}{HTML}{F7F7F7}
\definecolor{DarkGreen}{HTML}{1B6B3A}
\definecolor{AccentOrange}{HTML}{D97706}

\hypersetup{colorlinks=true, linkcolor=DeepNavy, urlcolor=DeepNavy}
\pagestyle{fancy}
\fancyhf{}
\lhead{Data Analysis Lesson}
\chead{Year-over-Year Sales Growth}
\rhead{Python + Power BI}
\cfoot{\thepage}
\renewcommand{\headrulewidth}{0.4pt}

\titleformat{\section}{\Large\bfseries\color{DeepNavy}}{}{0em}{}
\titleformat{\subsection}{\large\bfseries\color{DeepNavy}}{}{0em}{}

\lstdefinestyle{pythonstyle}{
    language=Python,
    backgroundcolor=\color{CodeBg},
    basicstyle=\ttfamily\small,
    keywordstyle=\color{DeepNavy}\bfseries,
    commentstyle=\color{DarkGreen},
    stringstyle=\color{AccentOrange},
    showstringspaces=false,
    breaklines=true,
    frame=single,
    rulecolor=\color{gray!40},
    tabsize=4
}

\lstdefinelanguage{DAX}{
  morekeywords={CALCULATE,SUM,DIVIDE,SAMEPERIODLASTYEAR,CALENDAR,MIN,MAX,YEAR,MONTH,FORMAT,IF,BLANK},
  sensitive=false,
  morecomment=[l]{//},
  morestring=[b]"
}
\lstdefinestyle{daxstyle}{
    language=DAX,
    backgroundcolor=\color{CodeBg},
    basicstyle=\ttfamily\small,
    keywordstyle=\color{DeepNavy}\bfseries,
    commentstyle=\color{DarkGreen},
    stringstyle=\color{AccentOrange},
    showstringspaces=false,
    breaklines=true,
    frame=single,
    rulecolor=\color{gray!40},
    tabsize=4
}

\setlength{\parindent}{0pt}
\setlength{\parskip}{5pt}

\begin{document}

\begin{center}
    {\Huge \bfseries \color{DeepNavy} Year-over-Year Sales Growth}\\[4pt]
    {\Large Using Python and Power BI DAX}\\[6pt]
    \textbf{Teacher: Nigussie Guluma} \quad | \quad Data Analysis Lesson
\end{center}

\begin{tcolorbox}[colback=LightBlue,colframe=DeepNavy,title=Learning Goal]
Students will understand how to calculate, interpret, and compare \textbf{Year-over-Year (YoY) Sales Growth} using two tools: \textbf{Python pandas} and \textbf{Power BI DAX}.
\end{tcolorbox}

\section{1. What is Year-over-Year Sales Growth?}
Year-over-Year Sales Growth compares sales from one year to the sales from the previous year. It helps answer this question:

\begin{center}
\textbf{``Compared to last year, did sales increase or decrease?''}
\end{center}

\begin{tcolorbox}[colback=LightGray,colframe=DeepNavy,title=Formula]
\[
\text{YoY Growth \%} = \frac{\text{Current Year Sales} - \text{Previous Year Sales}}{\text{Previous Year Sales}} \times 100
\]
\end{tcolorbox}

\textbf{Interpretation:}
\begin{itemize}[itemsep=2pt]
    \item Positive percentage means sales increased.
    \item Negative percentage means sales decreased.
    \item Zero percent means sales stayed the same.
\end{itemize}

\section{2. Sample Sales Data}
\begin{center}
\begin{tabular}{ccc}
\toprule
\textbf{Year} & \textbf{Sales} & \textbf{Meaning} \\
\midrule
2021 & \$40,000 & Starting year \\
2022 & \$46,000 & Sales increased \\
2023 & \$52,900 & Sales increased again \\
2024 & \$50,000 & Sales decreased \\
2025 & \$65,000 & Sales increased strongly \\
\bottomrule
\end{tabular}
\end{center}

\section{3. Manual Example}
Find the YoY growth from 2024 to 2025.

\[
\frac{65000 - 50000}{50000} \times 100
\]

\[
\frac{15000}{50000} \times 100 = 30\%
\]

\begin{tcolorbox}[colback=LightBlue,colframe=DeepNavy,title=Conclusion]
From 2024 to 2025, sales increased by \textbf{30\%}. This means the company sold much more in 2025 than in 2024.
\end{tcolorbox}

\section{4. Python Method Using pandas}
Python is useful when students want to calculate sales growth from a table or dataset. The pandas function \texttt{pct\_change()} calculates percent change from one row to the next row.

\subsection{Python Code}
\begin{lstlisting}[style=pythonstyle]
import pandas as pd
import matplotlib.pyplot as plt

# Step 1: Create the dataset
data = {
    "Year": [2021, 2022, 2023, 2024, 2025],
    "Sales": [40000, 46000, 52900, 50000, 65000]
}

df = pd.DataFrame(data)

# Step 2: Calculate previous year sales
df["Previous Year Sales"] = df["Sales"].shift(1)

# Step 3: Calculate YoY growth amount
df["YoY Growth"] = df["Sales"] - df["Previous Year Sales"]

# Step 4: Calculate YoY growth percent
df["YoY Growth %"] = df["Sales"].pct_change() * 100

print(df)
\end{lstlisting}

\subsection{Expected Python Output}
\begin{center}
\begin{tabular}{ccccc}
\toprule
\textbf{Year} & \textbf{Sales} & \textbf{Previous Year Sales} & \textbf{YoY Growth} & \textbf{YoY Growth \%} \\
\midrule
2021 & 40,000 & -- & -- & -- \\
2022 & 46,000 & 40,000 & 6,000 & 15.00\% \\
2023 & 52,900 & 46,000 & 6,900 & 15.00\% \\
2024 & 50,000 & 52,900 & -2,900 & -5.48\% \\
2025 & 65,000 & 50,000 & 15,000 & 30.00\% \\
\bottomrule
\end{tabular}
\end{center}

\subsection{Python Graph Code}
\begin{lstlisting}[style=pythonstyle]
# Bar graph for total sales by year
plt.figure(figsize=(8, 5))
plt.bar(df["Year"], df["Sales"])
plt.title("Total Sales by Year")
plt.xlabel("Year")
plt.ylabel("Sales")
plt.show()

# Line graph for YoY Growth Percentage
plt.figure(figsize=(8, 5))
plt.plot(df["Year"], df["YoY Growth %"], marker="o")
plt.title("Year-over-Year Sales Growth %")
plt.xlabel("Year")
plt.ylabel("YoY Growth %")
plt.axhline(0, linestyle="--")
plt.show()
\end{lstlisting}

\begin{tcolorbox}[colback=LightGray,colframe=DeepNavy,title=Python Explanation]
\texttt{shift(1)} moves the sales values down one row so each year can be compared with the previous year. \texttt{pct\_change()} calculates the percentage change automatically.
\end{tcolorbox}

\newpage
\section{5. Power BI Method Using DAX}
Power BI is useful when students want to build dashboards and interactive reports. DAX measures calculate values dynamically based on filters, slicers, and visuals.

\subsection{Data Model Requirement}
For YoY DAX calculations to work correctly, use a proper Date Table.

\begin{center}
\begin{tabular}{ll}
\toprule
\textbf{Table} & \textbf{Column} \\
\midrule
Sales & Sales[SalesAmount] \\
Sales & Sales[Date] \\
Calendar & Calendar[Date] \\
\bottomrule
\end{tabular}
\end{center}

Create a relationship:
\begin{center}
\textbf{Calendar[Date]  $\rightarrow$  Sales[Date]}
\end{center}

\subsection{Create a Date Table in DAX}
\begin{lstlisting}[style=daxstyle]
Calendar =
CALENDAR(
    MIN(Sales[Date]),
    MAX(Sales[Date])
)
\end{lstlisting}

Add useful columns:
\begin{lstlisting}[style=daxstyle]
Year = YEAR(Calendar[Date])
Month = FORMAT(Calendar[Date], "MMMM")
Month Number = MONTH(Calendar[Date])
\end{lstlisting}

\subsection{DAX Measures}
\textbf{Measure 1: Total Sales}
\begin{lstlisting}[style=daxstyle]
Total Sales =
SUM(Sales[SalesAmount])
\end{lstlisting}

\textbf{Measure 2: Previous Year Sales}
\begin{lstlisting}[style=daxstyle]
Sales Previous Year =
CALCULATE(
    [Total Sales],
    SAMEPERIODLASTYEAR(Calendar[Date])
)
\end{lstlisting}

\textbf{Measure 3: YoY Sales Growth Amount}
\begin{lstlisting}[style=daxstyle]
YoY Sales Growth =
[Total Sales] - [Sales Previous Year]
\end{lstlisting}

\textbf{Measure 4: YoY Sales Growth Percentage}
\begin{lstlisting}[style=daxstyle]
YoY Sales Growth % =
DIVIDE(
    [YoY Sales Growth],
    [Sales Previous Year],
    0
)
\end{lstlisting}

Format \textbf{YoY Sales Growth \%} as a percentage in Power BI.

\section{6. Python vs. Power BI Comparison}
\begin{center}
\begin{tabular}{p{0.3\textwidth}p{0.3\textwidth}p{0.3\textwidth}}
\toprule
\textbf{Question} & \textbf{Python} & \textbf{Power BI} \\
\midrule
Best use & Data analysis and coding practice & Dashboard and business reporting \\
Main tool & pandas & DAX measures \\
Formula method & \texttt{pct\_change()} or manual formula & \texttt{SAMEPERIODLASTYEAR()} and \texttt{DIVIDE()} \\
Output & DataFrame, graphs, calculations & Interactive tables, charts, slicers \\
Good for students because & They see how the calculation is built step by step & They learn business intelligence reporting \\
\bottomrule
\end{tabular}
\end{center}

\section{7. Student Practice}
Use the following data to answer the questions.

\begin{center}
\begin{tabular}{cc}
\toprule
\textbf{Year} & \textbf{Sales} \\
\midrule
2020 & \$35,000 \\
2021 & \$40,000 \\
2022 & \$46,000 \\
2023 & \$52,900 \\
2024 & \$50,000 \\
2025 & \$65,000 \\
\bottomrule
\end{tabular}
\end{center}

\begin{enumerate}[itemsep=4pt]
    \item Calculate the YoY growth amount from 2020 to 2021.
    \item Calculate the YoY growth percentage from 2021 to 2022.
    \item Which year had negative growth?
    \item Which year had the highest growth percentage?
    \item In Python, which function can calculate percent change automatically?
    \item In DAX, which function returns sales from the previous year?
    \item Explain in one sentence why YoY growth is useful for a business.
\end{enumerate}

\section{8. Exit Ticket}
\begin{enumerate}[itemsep=4pt]
    \item What does YoY stand for?
    \item Write the formula for YoY Growth \%.
    \item If sales increased from \$80,000 to \$100,000, what is the YoY Growth \%?
    \item Which tool is better for dashboards: Python or Power BI?
    \item Which tool is better for writing code and practicing data analysis: Python or Power BI?
\end{enumerate}

\begin{tcolorbox}[colback=LightBlue,colframe=DeepNavy,title=Teacher Note]
Both Python and Power BI use the same mathematical idea. Python helps students understand the calculation through code. Power BI helps students communicate results through dashboards and business reports.
\end{tcolorbox}

\end{document}
